\noindent\textbf{Resultado.} \(x(t)=H(t-1)[1-\cos(t-1)]-H(t-5)[1-\cos(t-5)]\).

\noindent\textit{Desarrollo.} La entrada es \(H(t-1)-H(t-5)\). Con condiciones nulas,
\[X=\frac{e^{-s}-e^{-5s}}{s(s^2+1)}.\]
La identidad \(1/[s(s^2+1)]=1/s-s/(s^2+1)\) corresponde a \(g(t)=1-\cos t\). Trasladar \(g\) a cada evento da la solución.
Antes de uno, \(x=0\); entre uno y cinco, \(x=1-\cos(t-1)\); después de cinco, \(x=\cos(t-5)-\cos(t-1)\). Cada rama satisface la ecuación con su entrada. Como \(g(0)=g'(0)=0\), posición y velocidad son continuas en ambos cambios, mientras que la aceleración tiene el salto impuesto por la fuerza. Esta es la solución \(C^1\), \(C^2\) por tramos; no se exige una segunda derivada continua en los cambios de fuerza.
Al apagarse la fuerza queda una oscilación libre acotada. No se deduce amortiguamiento de un modelo que no lo contiene. Laplace converge, por ejemplo, en \(\Re s>0\). El modelo se ha dado normalizado; no se inventa una conversión de sus unidades a un dispositivo físico particular.
